\chapter{Literature Review}

\section{The Evolution of Employee Retention Strategies}
Employee retention has historically been viewed through the lens of human psychology and organizational behavior. Early models of employee turnover, developed in the mid-20th century, focused heavily on the concepts of job satisfaction and perceived organizational commitment. Traditional retention strategies were largely uniform, offering standardized benefits, periodic salary increments, and generalized career development programs. While these broad-stroke strategies form the foundation of good HR practice, they suffer from a fundamental flaw: they are predominantly reactive and generalized.

As highlighted by Frye and Boom \cite{frye2018employee}, traditional retention policies are often activated only after an employee signals their intent to leave, usually by securing another offer or formally resigning. At this juncture, HR departments frequently resort to counter-offers. However, research indicates that employees who accept counter-offers tend to leave the organization shortly thereafter anyway, as the underlying drivers of their dissatisfaction—be it management style, lack of growth opportunities, or poor work-life balance—remain unaddressed. This reactive posture necessitates a paradigm shift towards identifying the antecedents of turnover long before the breaking point is reached.

\section{The Rise of People Analytics}
The transition from reactive HR to proactive talent management has been catalyzed by the advent of "People Analytics" (or HR Analytics). People analytics involves the application of statistics, mathematics, and predictive modeling to employee data to drive better business outcomes. Historically, HR departments collected vast amounts of data—performance reviews, compensation history, attendance records, and demographic information—but utilized it primarily for compliance and basic reporting.

The modern imperative, however, is to leverage this data strategically. By mining historical data, organizations can identify patterns and correlations that precede employee departures. The shift towards evidence-based HR allows decision-makers to move beyond intuition and anecdotal evidence. It transforms HR from an administrative function into a strategic partner capable of forecasting talent risks and optimizing workforce planning.

\section{Machine Learning Applications in HR}
The application of machine learning to predict employee attrition represents the cutting edge of people analytics. Unlike traditional statistical methods, which often assume linear relationships and struggle with complex, high-dimensional datasets, machine learning algorithms excel at uncovering non-linear, intricate patterns within human behavioral data.

Smith and Doe \cite{smith2018predicting} explored the efficacy of various machine learning techniques in predicting employee turnover. Their research demonstrated that while linear models like Logistic Regression provide excellent interpretability—allowing HR professionals to clearly understand the weights of different factors—non-linear models often achieve higher predictive accuracy. 

Similarly, Zhao and Hryniewicki \cite{zhao2019attrition} conducted a comparative study of machine learning models for attrition prediction. They emphasized the superiority of ensemble methods. Techniques such as Random Forest \cite{breiman2001random}, which construct a multitude of decision trees during training, and gradient boosting algorithms like XGBoost \cite{chen2016xgboost}, consistently outperform simpler models in capturing complex interactions between variables. For instance, while a low salary alone might not trigger attrition, a low salary combined with high overtime and a long commute might create a high-probability flight risk. Tree-based ensembles are particularly adept at identifying these multifaceted conditional relationships.

\section{Addressing Class Imbalance in Predictive Modeling}
A pervasive challenge in applying machine learning to employee attrition is the inherent class imbalance of the data. In any stable organization, the number of employees who remain (the majority class) vastly outnumbers those who leave (the minority class) within a given timeframe. If a dataset has an attrition rate of 15\%, a naive classifier that simply predicts "No Attrition" for every employee will boast an 85\% accuracy rate, despite failing entirely at its intended purpose of identifying flight risks.

This necessitates specialized techniques during the model training phase. One of the most prominent solutions in the literature is the Synthetic Minority Over-sampling Technique (SMOTE), introduced by Chawla et al. \cite{chawla2002smote}. Rather than merely duplicating existing minority instances (which can lead to severe overfitting), SMOTE synthesizes new, artificial minority examples by interpolating between existing minority instances in the feature space. This approach effectively balances the training distribution, compelling the learning algorithm to establish broader, more robust decision boundaries for the minority class, thereby significantly improving the recall of actual flight risks.

\section{Limitations of Existing Work and the Kognex Proposition}
While existing literature establishes a strong foundation for predictive HR analytics, several practical implementations fall short in real-world application. Many models prioritize raw accuracy over actionable metrics like Precision and Recall, leading to systems that either miss critical flight risks or generate too many false alarms, thereby inducing "alert fatigue" in HR teams. Furthermore, many academic models remain confined to Jupyter notebooks, lacking the intuitive user interfaces required for adoption by non-technical HR personnel.

The \textit{Kognex} project seeks to address these gaps. It not only leverages state-of-the-art algorithms (XGBoost) and advanced preprocessing (SMOTE) but also places a heavy emphasis on appropriate evaluation metrics tailored for imbalance. Crucially, Kognex transcends the theoretical by prioritizing deployment. By wrapping the predictive engine in a polished, accessible Streamlit application with a premium, focused user interface, Kognex bridges the gap between complex data science and actionable, everyday HR strategy.
